\chapter{成稿体裁与两份文件}
\label{ch:deliver}

本章规定成稿的篇幅、论证顺序和文件形态。大纲本身不是成稿。

\section{250 页所装的内容}

成稿正文不少于 250 页，按 A4、含图注与公式计。页数来自下列已经定义的对象，不来自重复段落。

\begin{center}
\begin{tabularx}{\textwidth}{lXr}
\toprule
部分 & 写满的对象 & 页 \\
\midrule
问题与闭环 & 定义、决策、隔离条件 & 16 \\
文献分层 & 第 \ref{ch:lit} 章每条来源的原文方程与边界 & 32 \\
架构 & 层、时序、合同、失效 & 18 \\
多尺度模型 & 推导、无量纲、闭合关系、参数区间 & 40 \\
估计 & 观测、病态、秩、集合更新 & 28 \\
对偶与探针 & 目标、探针库、损伤上界、信息上限 & 26 \\
因果 & 结构、三条路径、三类计算 & 22 \\
演化 & 算子、损失、层次更新、回退 & 20 \\
控制与安全 & 代价、约束、投影、降级、硬回路 & 24 \\
试验协议 & 因素、指标、软件边界 & 18 \\
命题证明 & 六条命题的完整证明 & 16 \\
图件 & 88 幅图及图注 & 42 \\
结果 & 仅日志中存在的 D 图 & 28 \\
实施例 & 三组步骤与 E 图 & 16 \\
\midrule
合计 &  & 约 342 \\
\bottomrule
\end{tabularx}
\end{center}

结果部分若日志未完成，正文保留协议与图位，不把空图位算作结果页，总页数仍由其余部分超过 250。结果页补上之后不删证明与文献。

\section{论证顺序}

每一章内部固定为四段，避免先宣布效果再补理由。

\begin{enumerate}
\item 本章要判定的量。
\item 继承的文献层与不继承的层。
\item 方程、命题或协议。
\item 该章不能推出的结论。
\end{enumerate}

摘要中的每句定量陈述必须指向一幅 D 图。没有 D 图的陈述只能出现在问题定义或命题中，并使用条件句。

\section{同一正文的两份文件}

成稿同时提交 PDF 与 Word。两份文件的章、节、公式编号、图号、表号和参考文献编号一致。

\begin{itemize}
\item PDF 由本书的 TeX 源生成，作为公式与图号的基准。
\item Word 由同一正文生成，用于批注和交付。图为独立图片文件插入，图号与 PDF 相同。
\item 若转换造成断行或公式外观差异，以 PDF 的公式编号为准，在 Word 中保持同一编号，不另起一套。
\item 两份文件使用同一封面、同一参考文献表和同一版日志索引。不出现版本号。
\item 演示程序、轨迹字段和参数字段以第 \ref{ch:demo} 章为准，不在两份文件里各写一套接口。
\end{itemize}

Word 不作为第二套技术内容的编辑处。技术修改先进入 TeX 源，再同步到 Word。

\section{演示所要求的计算形态}

交互演示要完整展示第 \ref{ch:demo} 章的四块与三条命题核对，计算不能只停留在浏览器里的显示脚本。

\begin{itemize}
\item 真值积分、集合更新、探针评分和投影放在可测试的数值核中。数值核使用第 \ref{ch:proof} 章的方程，输入输出为第 \ref{ch:demo} 章的字段。
\item 数值核要有独立于界面的测试：命题 \ref{prop:steady}、\ref{prop:pulse}、\ref{prop:budget}、\ref{prop:shield} 各至少一条可重复断言。
\item 界面读取数值核的状态流，展示轨迹、路径和投影。界面不另写一套差分。
\item 下载文件由数值核直接序列化。
\item 多孔介质离线结果以同一列名接入，界面不改判据。
\end{itemize}

在数值核和测试断言尚未存在时，成稿只保留本章与第 \ref{ch:demo} 章的规格，不附截图充当演示结果。
